% ------------------------------------------------------------------------
% Vortrag: Rennes 08.07.09
% ------------------------------------------------------------------------
% name: ruckd_spangl_rennes09.tex
% ------------------------------------------------------------------------
%
%\def\PRFourSided{1} % angeschaltet: 4seitig / auskommentiert: 1seitig
\def\PRStyles{1}
\def\PRSweave{1}
\def\PRBeamer{1}   % angeschaltet: Handout / auskommentiert: Vortrag
\def\PRfrutiger{1} % angeschaltet: Frutiger Fonts / auskommentiert: Times
%
%---------------------------------------------------------------------------
%
\def\LRGRA{\mbox{\parbox[t]{0.51cm}{\includegraphics[%bb=0 20 100 50, scale=.4,
width=.9cm]{Rlogo}}}}
\def\RGRA{\mbox{\parbox[t]{0.51cm}{\includegraphics[%bb=0 20 100 50, scale=.4,
width=.5cm]{Rlogo}}}}
%
%%%%%%%%%%%%%%%%%%%%%%%%%%%%%%%%%%%%%%%%%%%%%%%%%%%%%%%%%%%%%%%%%%%%%%%%%%%
\ifx\PRBeamer\undefined%
%
% we use handout
%
  \documentclass[handout]{beamer}%
%
\def\hup{0cm}
%\def\RGRA{\sf R}
%\def\LRGRA{\sf\LARGE R}
\def\nL{\newline}
%
\else%
%
% we use beamer slides
%
   \documentclass[10pt]{beamer}
\def\hup{-18cm}

\def\nL{\relax }
%
\fi
%%%%%%%%%%%%%%%%%%%%%%%%%%%%%%%%%%%%%%%%%%%%%%%%%%%%%%%%%%%%%%%%%%%%%%%%%%%
%
%Preamble:
%%
\usepackage{beamer_pre_ITWM}
\usepackage{RColorBrew}
%
%%%%%%%%%%%%%%%%%%%%%%%%%%%%%%%%%%%%%%%%%%%%%%%%%%%%%%%%%%%%%%%%%%%%%%%%%%%
%
% particular packages for this talk:
%
\RequirePackage{multirow}
\newcommand{\Person}[1]{\protect{\color{Dark2.3}\bf #1}}
\newcommand{\hreft}[1]{{\href{#1}{{\color{Dark2.4}\url{#1}}}}}
\newcommand{\pkg}[1]{\protect{\color{Dark2.3}\tt #1}}
\newcommand{\UrL}[1]{{\color{Dark2.4}\url{#1}}}
%
%%%%%%%%%%%%%%%%%%%%%%%%%%%%%%%%%%%%%%%%%%%%%%%%%%%%%%%%%%%%%%%%%%%%%%%%%%%
%
\title[Robustness for Kalman filtering revisited]%
{{\LARGE\LRGRA $\;\;$--Package~\pkg{robKalman} --- \newline}\normalsize \RGRA . Kalman's revenge
or \\ \RGRA obustness for Kalman Filtering \RGRA evisited}
%
\pdfinfo{ /title({R--Package 'robKalman' --- R. Kalman's revenge
or Robustness for Kalman Filtering revisited})
/author({Peter Ruckdeschel}) /keywords({Robustness, R, S4 classes, Kalman Filtering}) }
%
\author[Peter Ruckdeschel and Bernhard Spangl]{Peter Ruckdeschel\inst{1}
\and Bernhard Spangl\inst{2}\\[2ex]
%\protect{%
%{\parbox[b]{12cm}\begin{group}
\tiny $^1$ 
{%{\includegraphics[bb=262 426 291 472,width=0.35cm]{ITWMLOGO}}\hspace{0.3cm}
\tiny Fraunhofer ITWM, Kaiserslautern, Germany, \hspace{0.18cm}
\href{mailto:Peter.Ruckdeschel@itwm.fraunhofer.de}{\tt\tiny Peter.Ruckdeschel@itwm.fraunhofer.de}}\\
$^2$ %
{%{\includegraphics[bb=262 426 291 472,width=0.35cm]{RLI_Logo_klein_CMYK}}\hspace{0.3cm}
\tiny Universit�t f�r Bodenkultur, Vienna, Austria,\hspace{0.2cm}
\href{mailto:Bernhard.Spangl@boku.ac.at}{\tt\tiny Bernhard.Spangl@boku.ac.at}}
%\end{group}
%}}
}
%
%\institute%
%\href{http://www.rali.boku.ac.at/statedv.html}{\tiny \tt www.rali.boku.ac.at/statedv.html}}}
\lstset{language=R,basicstyle=\small\color{Rcolor}}
\date{Rennes, July 9, 2009}
%
%\LogoOff
%
\begin{document}
%\gdef\RGRA0{\RGRA}
%%%%%%%%%%%%%%%%%%%%%%%%%%%%%%%%%%%%%%%%%%%%%%%%%%%%%%%%%%%%%%%%%%%%%%%%%%%%%%%%%%%%%%%%%%%%%%%%%%%%%%%%%%%
%
% Folie 0
%
%%%%%%%%%%%%%%%%%%%%%%%%%%%%%%%%%%%%%%%%%%%%%%%%%%%%%%%%%%%%%%%%%%%%%%%%%%%%%%%%%%%%%%%%%%%%%%%%%%%%%%%%%%%
\frame[plain]{\titlepage}
\ifx\PRBeamer\undefined%
\relax%
\else
\begin{frame}
\makeextratitle
\end{frame}
\fi

%-------------------------------------------------------------------------------
\section{General setup}
%-------------------------------------------------------------------------------
\subsection{Euclidean State Space Models}
%-------------------------------------------------------------------------------
%%%%%%%%%%%%%%%%%%%%%%%%%%%%%%%%%%%%%%%%%%%%%%%%%%%%%%%%%%%%%%%%%%%%%%%%%%%%%%%%%%%%%%%%%%%%%%%%%%%%%%%%%%%
%
% Folie I-1
%
%%%%%%%%%%%%%%%%%%%%%%%%%%%%%%%%%%%%%%%%%%%%%%%%%%%%%%%%%%%%%%%%%%%%%%%%%%%%%%%%%%%%%%%%%%%%%%%%%%%%%%%%%%%
\begin{frame}{}
\frametitle{Euclidean State Space Models}
\structure{Definitions and Assumptions:}\newline
---
Time--Discrete, Euclidean Setup
%
\color{black} ideal model:
\begin{align}
x_t&=F(  x_{t-1},t)+ v_t, &%\quad
  v_t &\stackrel{\rm \SSs indep.}{\sim} (0,Q_t),&& [\mbox{$p$--dim}], \nonumber \\
y_t&=Z(  x_t,t)+ \ve_t,& %\quad
\ve_t&\stackrel{\rm \SSs indep.}{\sim} (0,V_t), &&[\mbox{$q$--dim}], \nonumber \\
&&\qquad x_0 &\stackrel{\phantom{\rm \SSs indep.}}{\sim}  (a_0,Q_0),&&[\mbox{$p$--dim}],\nonumber \\
\noalign{\qquad \mbox{$ \{v_t\}, \{\ve_t\},  x_0$ indep.\ as processes}}\nonumber
\end{align}
%
\uncover<2->{functions $F,Z$ smooth with known derivatives;\\
 hyper--parameters $Q_t,V_t,a_0$ known\\[2ex]}
\uncover<3->{
extensible to:
\begin{itemize}
  \item continuous time (SDE's)
  \item incorporate user-specified controls
\end{itemize}} 
\end{frame}
%

%-------------------------------------------------------------------------------
\subsection{Types of Outliers}
%-------------------------------------------------------------------------------
%%%%%%%%%%%%%%%%%%%%%%%%%%%%%%%%%%%%%%%%%%%%%%%%%%%%%%%%%%%%%%%%%%%%%%%%%%%%%%%%%%%%%%%%%%%%%%%%%%%%%%%%%%%
%
% Folie I-3
%
%%%%%%%%%%%%%%%%%%%%%%%%%%%%%%%%%%%%%%%%%%%%%%%%%%%%%%%%%%%%%%%%%%%%%%%%%%%%%%%%%%%%%%%%%%%%%%%%%%%%%%%%%%%
\begin{frame}{}
\frametitle{Types of Outliers}\vspace{-4ex}
%
\protect{\color{black}{\color{Dark2.2}
\begin{eqnarray}
\noalign{\mbox{\color{black}exogenous outliers affecting only singular observations}}\nonumber\\[-1.5ex]
%{\rm AO} &\color{black}::& \color{black}\ve^{\SSs\rm re}_t \sim (1-r_{\SSs\rm AO}){\cal N}_q(0,V_t)+\color{Dark2.2} r_{\SSs\rm AO} {\cal
%L}(\ve^{\SSs\rm cont}_t) \\
{\rm AO} &\color{black}::& \color{black}\ve^{\SSs\rm re}_t \sim
(1-r_{\SSs\rm AO}){\cal L}(\ve^{\SSs\rm id}_t)+
\color{Dark2.2} r_{\SSs\rm AO} {\cal L}(\ve^{\SSs\rm di}_t) \nonumber\\
{\rm SO} &\color{black}::& \color{black}y^{\SSs\rm re}_t  \sim (1-r_{\SSs\rm SO}){\cal
L}(y^{\SSs\rm id}_t)+\color{Dark2.2}r_{\SSs\rm SO}{\cal
L}(y^{\SSs\rm di}_t)
\nonumber\\[.5ex]
%{\rm IO} &\color{black}::& \color{black}v^{\SSs\rm real}_t  \sim (1-r_{\SSs\rm IO}){\cal N}_p(0,Q_t)+\color{Dark2.2}r_{\SSs\rm
%IO} {\cal L}(v^{\SSs\rm cont}_t)\color{black}
\noalign{\mbox{\color{black} endogenous outliers / structural changes}}\nonumber\\[-2.5ex]
{\rm IO} &\color{black}::& \color{black}v^{\SSs\rm re}_t  \sim
(1-r_{\SSs\rm IO}){\cal L}(v^{\SSs\rm id}_t)+
\color{Dark2.2}r_{\SSs\rm IO} {\cal L}(v^{\SSs\rm di}_t)\color{black}\nonumber
\end{eqnarray}
\color{black}}}
Notation: $\qquad$
%\begin{equation}
$\mybm{\,\,\cdot\,\,}\hat=\mybm{\,\,\cdot\,}^{\SSs\rm id},\qquad
 \hat{\mybm{\,\,\cdot\,\,}}\hat=\mybm{\,\,\cdot\,}^{\SSs\rm re},\qquad
 \color{Dark2.2}\tilde{\mybm{\,\,\cdot\,\,}}\hat=\mybm{\,\,\cdot\,}^{\SSs\rm di}\color{black}$
%\end{equation}
\vspace{1ex}

\uncover<2->{\structure{\large Different and competing goals}\vspace{.4ex}
\begin{itemize}
  \item<2->[A/SO\hphantom{ \& IO}] attenuation of ``false alarms''
  \item<2->[IO \hphantom{ \& A/SO}] $\!$tracking: detect structural changes as fast as possible;\\
  $\hspace{0.8cm}\!$ recovering: clean data from structural changes
  \item<2->[A/SO \& IO] identification problem:\\ $\hspace{0.8cm}$simultaneous treatment only possible with delay
\end{itemize}}
\end{frame}

%%%%%%%%%%%%%%%%%%%%%%%%%%%%%%%%%%%%%%%%%%%%%%%%%%%%%%%%%%%%%%%%%%%%%%%%%%%%%%%%%%%%%%%%%%%%%%%%%%%%%%%%%%%
%
% Folie I-4
%
%%%%%%%%%%%%%%%%%%%%%%%%%%%%%%%%%%%%%%%%%%%%%%%%%%%%%%%%%%%%%%%%%%%%%%%%%%%%%%%%%%%%%%%%%%%%%%%%%%%%%%%%%%%
%
%-------------------------------------------------------------------------------
\subsection{Classical Method: Kalman--Filter}
%-------------------------------------------------------------------------------
%%%%%%%%%%%%%%%%%%%%%%%%%%%%%%%%%%%%%%%%%%%%%%%%%%%%%%%%%%%%%%%%%%%%%%%%%%%%%%%%%%%%%%%%%%%%%%%%%%%%%%%%%%%
%
% Folie I-5
%
%%%%%%%%%%%%%%%%%%%%%%%%%%%%%%%%%%%%%%%%%%%%%%%%%%%%%%%%%%%%%%%%%%%%%%%%%%%%%%%%%%%%%%%%%%%%%%%%%%%%%%%%%%%
\begin{frame}{}
\frametitle{Classical Method: Kalman--Filter}
\structure{Filter Problem}
\begin{equation}
\mbox{\normalsize\color{Dark2.3} $\Ew \big| x_t-f_t(y_{1:t})\big|^2 = \min_{f_t}\,!,$ \color{black}} \nonumber
\end{equation}
$$\mbox{with  } y_{1:t}=(y_1,\ldots,y_t), \quad y_{1:0}:=\emptyset$$
General solution: {\normalsize\color{Dark2.3} $\Ew[ x_t|y_{1:t}]$\color{black}}
---difficult to compute \\[3ex]
\uncover<2->{\structure{Kalman--Filter} \hspace{1em}{\footnotesize assuming $F(x,t)=F_t x$, $Z(x,t)=Z_t x$}\\
optimal solution among linear filters --- {Kalman[/Bucy]~[60/61]}:
\begin{align}%[lrclrcl]
%
\!\!\!\mbox{\color{Dark2.3}\mbox{Initialization: }\color{black}}&\!\!\!\!&
 x_{0|0} &= a_0 \nonumber%\label{bet1}
 \\
\!\!\!\mbox{\color{Dark2.3}\mbox{Prediction: }\color{black}}&\!\!\!\!
& x_{t|t-1}&= F_t  x_{t-1|t-1} ,\nonumber%\label{bet2}%\\
\!\!\!\quad& \color{gray}[\Delta x_t&\color{gray}= x_t- x_{t|t-1}]%\nonumber
\\
\!\!\!\mbox{\color{Dark2.3}\mbox{Correction: }\color{black}}%\nonumber\\
&\!\!\!\!& x_{t|t}&=  x_{t|t-1} + M^0_t \Delta y_t, \nonumber% \label{bet3}%\\
\!\!\!\quad& \color{gray} [\Delta y_t &\color{gray}= y_t-Z_t x_{t|t-1}]%\nonumber
%
\end{align}
and corresponding recursions for the prediction/filtering error covariances
$\Sigma_{t|t[-1]}$ and the Kalman gain $M^0_t$
}
%
\end{frame}
%

%%%%%%%%%%%%%%%%%%%%%%%%%%%%%%%%%%%%%%%%%%%%%%%%%%%%%%%%%%%%%%%%%%%%%%%%%%%%%%%%%%%%%%%%%%%%%%%%%%%%%%%%%%%
%
% Folie I-7
%
%%%%%%%%%%%%%%%%%%%%%%%%%%%%%%%%%%%%%%%%%%%%%%%%%%%%%%%%%%%%%%%%%%%%%%%%%%%%%%%%%%%%%%%%%%%%%%%%%%%%%%%%%%%
%-------------------------------------------------------------------------------
%\subsection{Features of the Kalman--Filter}
%-------------------------------------------------------------------------------
\begin{frame}{}
\frametitle{Features of the Kalman--Filter}
\begin{itemize}
\item[$+$] an easy, understandable structure:\\
initialization, prediction, correction step
\item[$+$] correction step is easily evaluable and interpretable: it is linear !
\item[$+$] strict recursivity / Markovian structure:\\
all information from the past useful for the future is captured in the
value of $ x_{t|t-1}$.
\item[$-$] the correction step is linear and thus not robust, as $y$ enters unbounded;
\end{itemize}
\uncover<2->{Aim of robustification:
try to retain all {``$+$''}'s, revise {``$-$''}}
\end{frame}
%%

%\setcounter{frame}%{0}
%
% Folie II-1
%
%
%
% ------------------------------------------------------------------------
\section[Implementation: R-package robKalman]{Implementation: {\sf R}-package {\tt robKalman}}
% ------------------------------------------------------------------------
\subsection{Concept / Strategy}
% ------------------------------------------------------------------------
%%%%%%%%%%%%%%%%%%%%%%%%%%%%%%%%%%%%%%%%%%%%%%%%%%%%%%%%%%%%%%%%%%%%%%%%%%%%%%%%%%%%%%%%%%%%%%%%%%%%%%%%%%%
%
% Folie III-1
%
%%%%%%%%%%%%%%%%%%%%%%%%%%%%%%%%%%%%%%%%%%%%%%%%%%%%%%%%%%%%%%%%%%%%%%%%%%%%%%%%%%%%%%%%%%%%%%%%%%%%%%%%%%%
\begin{frame}%Frame 8
  \frametitle{{\sf R}-package {\tt  robKalman} --- Contents}
  \begin{itemize}
     \item<2-> Kalman filter: filter, Kalman gain, covariances
     \item<2-> ACM-filter: filter, multivariate version, GM-estimator
     \item<2-> rLS-filter: filter, calibration of clipping height
  \begin{itemize}
    \item<2-> AO/SO-robust version 
    \item<2-> IO-robust version
    \item<2-> with a certain delay joint treatment of AO/SO's \& IO's
  \end{itemize}%
     \item<3-> extensible to further recursive filters:\\
      $\leadsto$ general interface \lstinline{recursiveFilter}\\
      with arguments:
        \begin{itemize}
           \item data
           \item state space model (hyper parameters) \\ $\qquad${\color{Dark2.8}[will be: object of class SSM]}
           \item {} {\bf functions for the init./pred./corr.step}\\ $\qquad$
           {\color{Dark2.8}[will be: object containing them]}
           \item {} {\color{Dark2.8}[will be: control object]}
        \end{itemize}%
  \end{itemize}%
\end{frame}
% ------------------------------------------------------------------------
%%%%%%%%%%%%%%%%%%%%%%%%%%%%%%%%%%%%%%%%%%%%%%%%%%%%%%%%%%%%%%%%%%%%%%%%%%%%%%%%%%%%%%%%%%%%%%%%%%%%%%%%%%%
%
% Folie III-2
%
%%%%%%%%%%%%%%%%%%%%%%%%%%%%%%%%%%%%%%%%%%%%%%%%%%%%%%%%%%%%%%%%%%%%%%%%%%%%%%%%%%%%%%%%%%%%%%%%%%%%%%%%%%%
\begin{frame}%Frame 9
  \frametitle{Implementation concept}
\begin{itemize}
  \item<1-> \structure{Programming language}
  \begin{itemize}
    \item completely in {\tt S}, $\quad$
    %\item
    perhaps some code in {\tt C} later ($\leadsto$ \pkg{FKF})
  \end{itemize}%
  \item<2-> \structure{Use existing infrastructure: \pkg{zoo}, \pkg{timeSeries}}
  \begin{itemize}
%      \item time series classes: {\tt ts}, {\tt its}, {\tt irts}, {\tt zoo}, {\tt zoo.reg}, {\tt tframe}
      \item for: graphics, diagnostics, management of date/time
   \end{itemize}
   \item<3-> \structure{Code in different layers}
   \begin{itemize}
     \item internal functions: no {\tt S4}-objects, no time stamps\newline
     {\footnotesize (helps bringing in code by ``non-{\tt S4}-people'')}
     \item user interface: {\tt S4}-objects, time stamps
   \end{itemize}%
   \item<4-> \structure{Use generating functions for encapsulation}
   \begin{itemize}
     \item without using structured arguments:
   \begin{itemize}
     \item too many arguments $\leadsto$ user looses track
     \item prone to name mis-matchings (positional, partial matching)
     \item bad alternative: fix defaults\ldots
   \end{itemize}%
     \item have generating functions to produce control objects
     \item control objects may be reused
   \end{itemize}%
\end{itemize}%
\end{frame}

% ------------------------------------------------------------------------
\subsection{Implementation so far  and current devel topics}
% ------------------------------------------------------------------------
%%%%%%%%%%%%%%%%%%%%%%%%%%%%%%%%%%%%%%%%%%%%%%%%%%%%%%%%%%%%%%%%%%%%%%%%%%%%%%%%%%%%%%%%%%%%%%%%%%%%%%%%%%%
%
% Folie III-3
%
%%%%%%%%%%%%%%%%%%%%%%%%%%%%%%%%%%%%%%%%%%%%%%%%%%%%%%%%%%%%%%%%%%%%%%%%%%%%%%%%%%%%%%%%%%%%%%%%%%%%%%%%%%%
\begin{frame}%Frame 11
  \frametitle{Implementation so far}
\structure{Interfaces so far}
\begin{itemize}
  \item<1-> preliminary, ``{\tt S4}-free'' interfaces
  \begin{itemize}
     \item Kalman filter (in our context) \lstinline{KalmanFilter}
     \item rLS: \lstinline{rLSFilter}($=$\lstinline{rLS.AO.Filter}), \\ 
     $\mbox{\hphantom{ rLS:}}$ \lstinline{rLS.IO.Filter}, \lstinline{rLS.IOAO.Filter}%\\
     \item<1-> ACM: \lstinline{ACMfilt}, \lstinline{ACMfilter}, \lstinline{mACMfilter}
     \item<2-> all realized as wrappers to \lstinline{recursiveFilter}
  \end{itemize}%
%\item<3-> all required external packages available from \href{http://cran.r-project.org}{\tt CRAN}
%{\tt methods}, {\tt graphics}, {\tt startupmsg}, {\tt dse1}, {\tt dse2}, {\tt MASS}, {\tt limma}, {\tt robustbase}
\item<3-> availability:  \pkg{robKalman} version {\tt 0.3} (incl.\ demos)\\%[-1.5ex]
$\quad$       \fbox{\hreft{http://r-forge.r-project.org/projects/robkalman/}}\vspace{0.7ex}
\end{itemize}%
%\end{frame}
% ------------------------------------------------------------------------
%\subsection{Next steps}
% ------------------------------------------------------------------------
%%%%%%%%%%%%%%%%%%%%%%%%%%%%%%%%%%%%%%%%%%%%%%%%%%%%%%%%%%%%%%%%%%%%%%%%%%%%%%%%%%%%%%%%%%%%%%%%%%%%%%%%%%%
%
% Folie III-4
%
%%%%%%%%%%%%%%%%%%%%%%%%%%%%%%%%%%%%%%%%%%%%%%%%%%%%%%%%%%%%%%%%%%%%%%%%%%%%%%%%%%%%%%%%%%%%%%%%%%%%%%%%%%%
%\begin{frame}%Frame 13
%  \frametitle{Next steps --- to do in Wien :-)}
\uncover<4->{\structure{Almost ready:}
\begin{itemize}
  \item<4-> {\tt  S4} classes: for SSM's; for output-classes; for method-classes;  \\
    \visible<beamer:0 | handout:0>{{\tt  S4} classes:} for control-classes (reuse \pkg{robustbase}-code)
  \item<5-> interfaces between {\tt S4}-layer and {\tt S4}-free layer\\
  %\item<5->
  \visible<beamer:0 | handout:0>{interfaces} to other SSM packages  \\
  \visible<beamer:0 | handout:0>{interfaces} to \pkg{robfilter} (\Person{Roland Fried} \& \Person{K. Schettlinger})
 \end{itemize}}%
\end{frame}
%
%%%%%%%%%%%%%%%%%%%%%%%%%%%%%%%%%%%%%%%%%%%%%%%%%%%%%%%%%%%%%%%%%%%%%%%%%%%%%%%%%%%%%%%%%%%%%%%%%%%%%%%%%%%
\section[Work in process]{Work in process}
%%%%%%%%%%%%%%%%%%%%%%%%%%%%%%%%%%%%%%%%%%%%%%%%%%%%%%%%%%%%%%%%%%%%%%%%%%%%%%%%%%%%%%%%%%%%%%%%%%%%%%%%%%%
%
% Folie IV-1
%
%%%%%%%%%%%%%%%%%%%%%%%%%%%%%%%%%%%%%%%%%%%%%%%%%%%%%%%%%%%%%%%%%%%%%%%%%%%%%%%%%%%%%%%%%%%%%%%%%%%%%%%%%%%
\begin{frame}{}
\frametitle{Work in process}
\structure{Release Plans}
\begin{itemize}
  \item package \pkg{robKalman} should be on \UrL{CRAN} by UseR! 2009, but\ldots
  \item at least: release on \UrL{CRAN} by end of August
  \item till then: refer to \UrL{r-forge}
\end{itemize}
\uncover<2->{\structure{Extensions}
\begin{itemize}
\item robust smoothing \hfill {\footnotesize\color{Dark2.8}($80\%$ done)}
\item robust EM-Algorithm to estimate unknown hyper parameters\\ (extending Shumway/Stoffer)
\hfill {\footnotesize\color{Dark2.8}($70\%$ done)}
\item interpretation as random coefficient regression \\$\leadsto$
robust regression-type approach ({\bf rIC}, {\bf mIC}) \hfill {\footnotesize\color{Dark2.8}($30\%$ done)}
\item connecttion to particle filters --- \\theory and computer interface  \hfill  {\footnotesize\color{Dark2.8}($10\%$ done)}
\item speeding up things / bridging to fast Kalman filter of \\ \pkg{FKF} by \Person{David Luethi}, \Person{Philipp Erb}  \hfill {\footnotesize\color{Dark2.8}($1\%$ done)}
\end{itemize}}
\end{frame}
% ------------------------------------------------------------------------
\section[{Collaboration on r-forge}]{Collaboration on {\tt r-forge}}
% ------------------------------------------------------------------------
\begin{frame}{}
\frametitle{Some experiences on collaborative programming on {\tt r-forge}}
\begin{itemize}
\item \structure{\UrL{r-forge}}:\newline {\bf very neat} for collaborative {\sf R} package development
\begin{footnotesize}
\begin{itemize}
\item version management ({\tt svn})
\item mail-forwarded {\tt log}-files of committed code \\$\leadsto$ keep track of work of others
\item bug tracker, archived mailing lists, \ldots
\item see slides by \Person{Stefan Theussl}
\end{itemize}
\end{footnotesize}
\item<2->needs serious conceptional preparations  
\begin{footnotesize}
\begin{itemize}
\item<2-> for separating/modularizing tasks
\item<2-> consistency: coding \& documentation conventions 
\end{itemize}
\end{footnotesize}
\item<2-> helpful: scheduling, reminders/deadlines for collaborators\ldots
\item<3-> summarizing:\\[.5ex] \structure{\textit{Collaborative programming is enjoyable and very exciting!}}
\end{itemize}\vspace{1ex}
\uncover<4->{\centerline{\structure{\Large \sc Thanks for your attention!}}}
\end{frame}
%%%%%%%%%%%%%%%%%%%%%%%%%%%%%%%%%%%%%%%%%%%%%%%%%%%%%%%%%%%%%%%%%%%%%%%%%%%%%%%%%%%%%%%%%%%%%%%%%%%%%%%%%%%
%
% Folie IV-2
%
%%%%%%%%%%%%%%%%%%%%%%%%%%%%%%%%%%%%%%%%%%%%%%%%%%%%%%%%%%%%%%%%%%%%%%%%%%%%%%%%%%%%%%%%%%%%%%%%%%%%%%%%%%%
%%%%%%%%%%%%%%%%%%%%%%%%%%%%%%%%%%%%%%%%%%%%%%%%%%%%%%%%%%%%%%%%%%%%%%%%%%%%%%%%%%%%%%%%%%%%%%%%%%%%%%%%%%%
%
% Folie IV-3
%
%%%%%%%%%%%%%%%%%%%%%%%%%%%%%%%%%%%%%%%%%%%%%%%%%%%%%%%%%%%%%%%%%%%%%%%%%%%%%%%%%%%%%%%%%%%%%%%%%%%%%%%%%%%

%\mode<handout>{
%\input{lit}
%\begin{frame}%<1-| handout:0>
%[allowframebreaks]%Frame 24
%  \frametitle<presentation>{Bibliography}
%\protect{\begin{tiny}
%\bibliography{litdb1}
%\end{tiny}}
%\end{frame}
%}
%
%\ifx\PRBeamer\undefined%
%% ------------------------------------------------------------------------
\section{References}
%% ------------------------------------------------------------------------
\begin{frame}%<presentation:0|handout>
[allowframebreaks]%[fragile,allowframebreaks]
\frametitle{References}\vspace{-1ex}%
\begin{scriptsize}%\vspace{-3ex}
\begin{description}\scriptsize
\item[Birmiwal, K. and Shen, J. (1993)]:
  \textit{Optimal robust filtering.} { Stat. Decis.\/}, {\bf 11}(2): 101--119.\vspace{-.4ex}
 \item[Durbin, J. and Koopman, S. J. (2001)]: \textit{Time Series Analysis by
   State Space Methods.} Oxford University Press.\vspace{-.4ex}
%  \item[Gilbert, P. (2005)]: Brief User's Guide: Dynamic Systems Estimation (DSE).
%{%\footnotesize
%Available in the file\newline {\tt doc/dse-guide.pdf} distributed together with the {\sf R} bundle
%  {\tt dse}, to be downloaded from \newline\hreft{http://cran.r-project.org}}
\item[Fried, R. and Schettlinger, K. (2008)]: {{\sf R}-package {\tt robfilter}: Robust Time Series Filters.}
\hreft{http://cran.r-project.org/web/packages/robfilter}.\vspace{-.4ex}
\item[Kalman, R.E. (1960)]: {A new approach to linear
  filtering and prediction problems.}\vspace{-.4ex}
 {\em Journal of Basic Engineering---Transactions of the ASME\/}, {\bf
  82}: 35--45.\vspace{-.4ex}
\item[Kalman, R.E. and Bucy, R. (1961)]:
  {New results in filtering and prediction theory.}
 {\em Journal of Basic Engineering---Transactions of the ASME\/}, {\bf
  83}: 95--108.\vspace{-.4ex}
  \item[Martin, D. (1979)]:
  \textit{Approximate conditional-mean type smoothers and interpolators.}
  In  \emph{Smoothing techniques for curve estimation.\/}
  Proc. Workshop Heidelberg 1979. Lect. Notes Math.~757, p.~117-143\vspace{-.4ex}
  \item[Masreliez C.J. and Martin R. (1977)]: {Robust Bayesian estimation for the linear model and
  robustifying the Kalman filter.}
 {\em IEEE Trans. Autom. Control\/}, {\bf {AC-22}}: 361--371.\vspace{-.4ex}
  \item[Ruckdeschel, P. (2001)]: \textit{Ans{\"a}tze zur Robustifizierung des Kalman Filters.}
  Bayreuther Mathematische Schriften, Vol. 64.
  \item[R Development Core Team (2009)]: \textit{{\sf R}: A language and environment for statistical computing\/}.
R Foundation for Statistical Computing, Vienna, Austria.\\
  \hreft{http://www.R-project.org}
  \item[R-Forge Administration and Development Team (2008)]: \textit{{R-Forge User's Manual, {\sc beta}\/}.}
{\tt SVN} revision: 47, August, 12 2008. \newline
  \hreft{http://r-forge.r-project.org/R-Forge_Manual.pdf}
\item [Schick, I.C. (1989)]: {\em {Robust recursive
  estimation of a discrete--time stochastic linear dynamic system in the
  presence of heavy-tailed observation noise}\/}.
 {Dissertation}, {Massachusetts Institute of Technology}, {Cambridge,
  MA}.
  \item[Schick I.C. and Mitter S.K. (1994)]:
  {Robust recursive estimation in the presence of heavy-tailed observation
  noise.} {\em Ann. Stat.\/}, {\bf 22}(2): 1045--1080.
  \item[Shumway, R.H. and Stoffer, D.S. (1982)]: \textit{An approach to time series smoothing and forecasting using the EM algorithm.}
  Journal of Time Series Analysis, \textbf{3}: 253--264.
  \item[Spangl, B. (2008)]: \textit{On Robust Spectral Density
Estimation\/}. PhD Thesis at Technical University, Vienna.\bigskip
\end{description}
\end{scriptsize}
\end{frame}
%\fi

\end{document}
